% Real-Space Boson Hamiltonian
% Converted on 2026-10-08 from docs/lswt/01-derivation/real-space-boson-hamiltonian.md (status: draft, source-section: Bosonic representations for spin Hamiltonian).
% Review status: draft; awaiting the user's physics and mathematics review.

\hypertarget{sec-lswt-real-space-boson-hamiltonian}{%
\section{Real-Space Boson Hamiltonian}\label{sec-lswt-real-space-boson-hamiltonian}}

The quadratic boson Hamiltonian collects the terms of second order in the Holstein--Primakoff operators. The expression below includes inter-site exchange and the Zeeman term in the local frame. The single-ion term with matrix \(\mathsf D_I\) is added in the subsection on the single-ion term below.

The circular-component expressions below are provisional: the contraction convention relating \(\widetilde J_\ell^{pq}\) to the Cartesian exchange matrix is not fully specified. The labels \(p,q\in\{+,-,0\}\) denote local circular components, rather than laboratory Cartesian indices.

\subsection{Quadratic Exchange and Field Terms}

For a bond \(\ell=(I,J)\), define the spin-weighted exchange coefficient

\[
t_\ell^{pq}=\sqrt{S_I S_J}\,\widetilde J_\ell^{pq},
\qquad \ell=(I,J).
\]

Retaining the quadratic terms gives the provisional component form

\[
\begin{aligned}
\hat H_2
=&\sum_{\ell=(I,J)\in\mathcal L}\Big[
t_\ell^{--}\hat a_I\hat a_J
+t_\ell^{-+}\hat a_I\hat a_J^\dagger
+t_\ell^{+-}\hat a_I^\dagger\hat a_J
+t_\ell^{++}\hat a_I^\dagger\hat a_J^\dagger
\\
&\hspace{7em}
-\widetilde J_\ell^{00}\left(S_I\hat n_J+S_J\hat n_I\right)
\Big]
+\sum_I\widetilde h_I^0\hat n_I.
\end{aligned}
\]

The first four terms contain the normal and anomalous boson products associated with each link. The remaining terms multiply local number operators. The field contribution involves the longitudinal component of the rotated Zeeman-energy vector, \(\widetilde h_I^0\), rather than an unrotated laboratory component.

\subsection{Local Number-Operator Coefficient}

Each site receives exchange contributions from all incident links. Let \(\partial I\) be the set of links incident on \(I\), and let \(I'_\ell\) be the opposite endpoint of link \(\ell\). For the exchange and field terms displayed above, the local coefficient is

\[
\mu_I
=\widetilde h_I^0
-\sum_{\ell\in\partial I}S_{I'_\ell}\widetilde J_\ell^{00}.
\]

This coefficient combines the rotated field with the longitudinal exchange contributions from neighboring spins. Its use inherits the circular-component convention and interaction scope stated above.

\subsection{Single-Ion Term}

The single-ion term \(\hat{\mathbf S}_I^{\mathsf T}\mathsf D_I\hat{\mathbf S}_I\) contains two spin operators of the same site, and their product is not determined by the classical vector alone. We use the rule that reproduces the spin-coherent expectation value of the classical energy in \lswtnoteref{LSWT spin Hamiltonian}{LSWT Overview}: the term enters the quadratic Hamiltonian as a self-coupling of site \(I\), expanded like a link whose two endpoints coincide, with the rotated matrix \(\widetilde{\mathsf D}_I=\mathsf R_I^{\mathsf T}\mathsf D_I\mathsf R_I\) replaced by
\[
\Big(1-\frac{1}{2S_I}\Big)\widetilde{\mathsf D}_I,
\]
and without the constant produced by normal-ordering the transverse operators of the same site. The corresponding classical energy is \(S_I(S_I-\tfrac12)\,\mathbf n_I^{\mathsf T}\mathsf D_I\mathbf n_I+\tfrac{S_I}{2}\operatorname{Tr}\mathsf D_I\); at \(S_I=\tfrac12\) it is the constant \(\operatorname{Tr}\mathsf D_I/4\) and the term drops out of the quadratic Hamiltonian, as the spin-\(\tfrac12\) identity of \lswtnoteref{LSWT spin Hamiltonian}{Bilinear Spin Hamiltonian} requires.

For an easy-axis term \(D(\hat S_I^z)^2\) with \(D<0\) and the spin along \(z\), only \(\widetilde D_I^{00}=D\) is nonzero, and the longitudinal part \((S_I-\hat n_I)^2\) contributes \(-2S_I(1-\tfrac1{2S_I})D\,\hat n_I=(2S_I-1)|D|\,\hat n_I\). The resulting single-ion gap \((2S_I-1)|D|\) equals the exact spacing between the levels \(m=S_I\) and \(m=S_I-1\), whereas the factor \(2S_I\) of a direct expansion with \(\widetilde{\mathsf D}_I\) overestimates it. The rule is justified by these two agreements, with the coherent-state classical energy and with the exact single-ion gap, rather than derived from the Holstein--Primakoff expansion of the operator product.

\subsection{References}

\subsubsection{Internal Documents}

\begin{itemize}
\tightlist
\item
  \lswtnoteref{LSWT spin Hamiltonian}{Bilinear Spin Hamiltonian}: defines the spin Hamiltonian expanded here.
\item
  \lswtnoteref{LSWT classical order}{Classical Order and Local Frame}: defines the rotated exchange tensor and field.
\item
  Holstein--Primakoff Expansion (Section~\ref{sec-lswt-holstein-primakoff-expansion}): supplies the boson expansion and harmonic truncation.
\item
  Momentum-Space BdG Hamiltonian (Section~\ref{sec-lswt-momentum-space-bdg-hamiltonian}): transforms this real-space quadratic Hamiltonian to momentum space.
\end{itemize}

\subsubsection{External Sources}

\begin{itemize}
\tightlist
\item
  None.
\end{itemize}
